% LTeX: enabled=false
% Sección 3: correspondencia entre objetivos y resultados, y próximos pasos

\seccion{Balance}

% ===============================================================
% 12. Objetivos y resultados
% ===============================================================
\newcommand{\etiq}[2]{\fcolorbox{#1}{#1!10}{\textcolor{#1}{\textbf{#2}}}}

\begin{frame}{Lo que nos propusimos y lo que logramos}
    \begin{tabla}{colspec = {Q[c,m,0.7cm] Q[l,m,4.3cm] X[l,m] Q[c,m,2.3cm]},
                  rows = {font=\footnotesize}, rowsep = 1pt}
        Obj. & Qué nos propusimos & Qué logramos & Estado \\
        1 & Analizar un conjunto de datos público & CornOrb seleccionado y revisado, con los datos a corregir identificados & \etiq{verde}{Alcanzado} \\
        2 & Construir el modelo de IA & Diseño definido y comparación de redes base programada, sin resultados todavía & \etiq{azul}{En proceso} \\
        3 & Integrar la explicación visual & Prevista con Grad-CAM, al cierre del modelado & \etiq{gray}{Por empezar} \\
        4 & Desarrollar la aplicación web & Partes y herramientas definidas (React y FastAPI) & \etiq{gray}{Por empezar} \\
        --- & Conseguir datos locales para validar & 9 centros contactados, charla con Roberto Albertazzi y trámites en curso & \etiq{azul}{En proceso} \\
    \end{tabla}
    \vspace{0.2cm}
    \centering
    {\small El avance coincide con el cronograma: la fase 1 está terminada y la fase 2 está en curso.}
\end{frame}

% ===============================================================
% 13. Próximos pasos
% ===============================================================
\begin{frame}{Próximos pasos}
    \begin{columns}[T]
        \begin{column}{0.48\textwidth}
            \begin{block}{Noviembre}
                \begin{itemize}
                    \item Corregir los datos según la revisión.
                    \item Entrenar y evaluar el modelo, priorizando no pasar por alto casos de queratocono.
                    \item Integrar el mapa de calor.
                    \item Empezar la aplicación web.
                \end{itemize}
            \end{block}
        \end{column}
        \begin{column}{0.48\textwidth}
            \begin{block}{Diciembre}
                \begin{itemize}
                    \item Terminar la aplicación web.
                    \item Conectar el modelo con la aplicación.
                    \item Tener el prototipo funcionando de punta a punta.
                \end{itemize}
            \end{block}
            \vspace{0.2cm}
            \begin{block}{En paralelo}
                \begin{itemize}
                    \item Avanzar con los trámites para los datos locales.
                \end{itemize}
            \end{block}
        \end{column}
    \end{columns}
    \idea{Enero queda como reserva por si algo se atrasa.}
\end{frame}
